\documentclass[11pt,a4paper]{article}

\usepackage[utf8]{inputenc}
\usepackage[T1]{fontenc}
\usepackage[portuguese]{babel}
\usepackage{geometry}
\usepackage{listings}
\usepackage{xcolor}
\usepackage{hyperref}
\usepackage{enumitem}

\geometry{margin=2.2cm}

\hypersetup{
    colorlinks=true,
    linkcolor=blue,
    urlcolor=blue
}

\definecolor{codegreen}{rgb}{0,0.6,0}
\definecolor{codegray}{rgb}{0.5,0.5,0.5}

\lstset{
    basicstyle=\ttfamily\small,
    breaklines=true,
    frame=single,
    numbers=left,
    numberstyle=\tiny\color{codegray},
    keywordstyle=\color{blue},
    commentstyle=\color{codegreen},
    showstringspaces=false
}

\setlist{noitemsep, topsep=2pt}

\title{
    Processamento de Linguagens \\
    \large Trabalho Prático \\
    \vspace{0.3cm}
    Compilador de Fortran 77 para Máquina Virtual
}

\author{
    Grupo 47 \\
    José Miguel Paredes Sampaio (A106908) \\
    José Pedro Joaquim Pereira (A106874) \\
    Sofia Margarida Freitas (A106798)
}

\date{maio de 2026}

\begin{document}

\maketitle
\tableofcontents
\newpage

\section{Resumo}

No âmbito da unidade curricular de Processamento de Linguagens, foi desenvolvido um compilador para um subconjunto de Fortran 77, com geração de código para a máquina virtual fornecida.

O compilador foi implementado em Python com recurso à biblioteca PLY (\textit{Python Lex-Yacc}) e encontra-se dividido em análise léxica, análise sintática, análise semântica, otimização simples da AST e geração de código. O projeto suporta variáveis escalares, arrays unidimensionais, expressões aritméticas, relacionais e lógicas, ciclos \texttt{DO}, condicionais \texttt{IF/THEN/ELSE/ENDIF}, \texttt{GOTO}, leitura e escrita, strings, subrotinas, funções e algumas otimizações locais.

As maiores dificuldades estiveram relacionadas com arrays, labels, controlo de fluxo e passagem de argumentos para subprogramas.

\section{Introdução}

O objetivo principal do projeto foi construir um compilador capaz de traduzir programas escritos num subconjunto de Fortran 77 para assembly da máquina virtual usada na unidade curricular.

Fortran 77 é uma linguagem clássica, muito usada historicamente em aplicações científicas e numéricas. Apesar da sua sintaxe relativamente simples, apresenta desafios relevantes para este trabalho, nomeadamente ciclos baseados em labels, arrays com indexação iniciada em 1 e subprogramas.

Além de implementar as construções base pedidas, o grupo tentou incluir algumas funcionalidades de valorização, nomeadamente funções, subrotinas, ciclos descendentes e uma fase simples de otimização. Estas extensões obrigaram a reorganizar parte da análise semântica e da geração de código, porque passaram a existir scopes diferentes e chamadas entre partes distintas do programa.

\section{Arquitetura do Compilador}

O projeto foi dividido nos ficheiros principais \texttt{lexer.py}, \texttt{parser.py}, \texttt{semantic\_analyzer.py}, \texttt{optimizer.py}, \texttt{translator.py} e \texttt{compiler.py}.

O fluxo geral é:

\begin{center}
Código Fonte $\rightarrow$ Lexer $\rightarrow$ Parser $\rightarrow$ Análise Semântica $\rightarrow$ AST $\rightarrow$ Otimizador $\rightarrow$ Código VM
\end{center}

O ficheiro \texttt{compiler.py} coordena o processo: lê o programa de entrada, invoca o parser, aplica otimizações simples e chama o tradutor para gerar o ficheiro final em código VM.

Cada módulo ficou responsável por uma parte concreta:

\begin{itemize}
    \item \texttt{lexer.py}: definição dos tokens e palavras reservadas;
    \item \texttt{parser.py}: regras gramaticais e construção da AST;
    \item \texttt{semantic\_analyzer.py}: validação de tipos, declarações, labels e subprogramas;
    \item \texttt{optimizer.py}: simplificações locais sobre a AST;
    \item \texttt{translator.py}: geração das instruções da máquina virtual.
\end{itemize}

\section{Análise Léxica}

A análise léxica foi implementada em \texttt{lexer.py} com \texttt{ply.lex}. Foram definidos tokens para palavras reservadas, identificadores, constantes, operadores e delimitadores.

As principais palavras reservadas reconhecidas são:

\begin{itemize}
    \item \texttt{PROGRAM}, \texttt{END}, \texttt{INTEGER}, \texttt{REAL}, \texttt{LOGICAL};
    \item \texttt{IF}, \texttt{THEN}, \texttt{ELSE}, \texttt{ENDIF};
    \item \texttt{DO}, \texttt{CONTINUE}, \texttt{GOTO};
    \item \texttt{READ}, \texttt{PRINT};
    \item \texttt{FUNCTION}, \texttt{SUBROUTINE}, \texttt{CALL}, \texttt{RETURN}.
\end{itemize}

Também são reconhecidos números inteiros, números reais, strings, operadores aritméticos, relacionais e lógicos. Os valores \texttt{.TRUE.} e \texttt{.FALSE.} são representados internamente como 1 e 0.

Os operadores relacionais suportados incluem a notação típica de Fortran, como \texttt{.LT.}, \texttt{.LE.}, \texttt{.GT.}, \texttt{.GE.}, \texttt{.EQ.} e \texttt{.NE.}. Também foram considerados operadores lógicos como \texttt{.AND.}, \texttt{.OR.} e \texttt{.NOT.}.

Exemplo de token:

\begin{lstlisting}[language=Python]
def t_NUMBER(t):
    r"\d+"
    t.value = int(t.value)
    return t
\end{lstlisting}

O compilador ignora espaços, tabulações e comentários iniciados por \texttt{!}. No pré-processamento também são ignorados alguns comentários clássicos de Fortran, com \texttt{C} ou \texttt{*} na primeira coluna.

\section{Análise Sintática}

A análise sintática foi implementada em \texttt{parser.py} com \texttt{ply.yacc}. A estrutura geral aceite pelo parser é, de forma simplificada:

\begin{lstlisting}
program : PROGRAM IDENTIFIER line_end declarations statements END opt_program_name subprograms
\end{lstlisting}

A gramática suporta declarações, atribuições, expressões, arrays, ciclos \texttt{DO}, condicionais, leitura, escrita, labels, \texttt{GOTO}, subrotinas e funções.

Foram também criadas regras para:

\begin{itemize}
    \item listas de variáveis em declarações;
    \item listas de argumentos em chamadas de subprogramas;
    \item expressões com precedência aritmética e lógica;
    \item blocos \texttt{IF} com ou sem \texttt{ELSE};
    \item ciclos \texttt{DO} com passo explícito ou implícito.
\end{itemize}

Para simplificar a distinção entre variáveis, arrays e chamadas de função, foi introduzido o conceito de \textit{designator}. Assim, uma construção como \texttt{A(I)} pode representar um acesso a array ou uma chamada de função, dependendo da informação semântica disponível.

Exemplo de regra de atribuição:

\begin{lstlisting}[language=Python]
def p_assignment(p):
    """
    assignment : designator EQ expression
    """
\end{lstlisting}

Foi necessário definir precedências entre operadores para evitar ambiguidades em expressões aritméticas, relacionais e lógicas.

Um aspeto importante foi permitir que os subprogramas fossem chamados antes da sua definição textual. Para isso, antes do parsing completo, o código é percorrido de forma simples para pré-declarar cabeçalhos de \texttt{FUNCTION} e \texttt{SUBROUTINE}.

\section{Análise Semântica}

A análise semântica foi implementada principalmente em \texttt{semantic\_analyzer.py}. Esta fase verifica:

\begin{itemize}
    \item uso de variáveis apenas depois de declaradas;
    \item declarações duplicadas;
    \item compatibilidade de tipos;
    \item acessos corretos a arrays;
    \item labels usados em \texttt{DO} e \texttt{GOTO};
    \item existência e assinatura de funções e subrotinas;
    \item número correto de argumentos em chamadas.
\end{itemize}

A tabela de símbolos guarda, para cada identificador, o tipo e a dimensão, caso seja um array. Os endereços de memória são atribuídos posteriormente pelo tradutor.

Com a introdução de subprogramas, a análise semântica passou também a gerir scopes. O programa principal tem a sua tabela de símbolos, e cada função ou subrotina possui uma tabela própria. Os parâmetros formais têm de ser declarados dentro do subprograma, permitindo ao compilador conhecer o seu tipo e perceber se representam escalares ou arrays.

Também são mantidas estruturas separadas para labels de cada scope, evitando que um \texttt{GOTO} dentro de uma subrotina seja validado contra labels do programa principal.

Nas funções, foi seguida a convenção de Fortran: o valor de retorno é definido por uma atribuição ao identificador com o mesmo nome da função. Por exemplo:

\begin{lstlisting}
INTEGER FUNCTION DOBRO(N)
INTEGER N
DOBRO = N * 2
RETURN
END
\end{lstlisting}

A coerção implícita suportada é limitada: valores \texttt{INTEGER} podem ser atribuídos a variáveis \texttt{REAL}, sendo posteriormente emitida a conversão adequada durante a geração de código.

\section{Funcionalidades Suportadas}

O subconjunto implementado suporta:

\begin{itemize}
    \item tipos \texttt{INTEGER}, \texttt{REAL} e \texttt{LOGICAL};
    \item variáveis escalares e arrays unidimensionais;
    \item expressões aritméticas, relacionais e lógicas;
    \item \texttt{READ *} e \texttt{PRINT *}, incluindo strings;
    \item condicionais \texttt{IF/THEN/ELSE/ENDIF};
    \item ciclos \texttt{DO} com label, \texttt{CONTINUE} e passo positivo ou negativo;
    \item \texttt{GOTO};
    \item \texttt{SUBROUTINE}, \texttt{CALL}, \texttt{FUNCTION} e \texttt{RETURN};
    \item passagem de argumentos escalares e arrays unidimensionais para subrotinas;
    \item chamadas de funções em expressões;
    \item otimizações simples antes da geração de código.
\end{itemize}

\section{Otimizações}

Foi implementado o módulo \texttt{optimizer.py}, que percorre a AST antes da geração de código. As otimizações aplicadas são locais e simples:

\begin{itemize}
    \item avaliação de expressões constantes;
    \item simplificações como \texttt{X + 0}, \texttt{X * 1} e \texttt{X * 0};
    \item simplificação de expressões lógicas constantes;
    \item remoção de atribuições do tipo \texttt{X = X};
    \item remoção de ramos condicionais com condição constante, quando não há labels que possam ser afetados.
\end{itemize}

Estas otimizações não pretendem substituir uma fase de otimização avançada, mas permitem reduzir algum código desnecessário e demonstrar uma etapa adicional comum em compiladores.

Por exemplo, uma expressão como:

\begin{lstlisting}
A = (2 + 3) * 4
\end{lstlisting}

pode ser simplificada antes da tradução para:

\begin{lstlisting}
A = 20
\end{lstlisting}

Também são removidas atribuições sem efeito, como \texttt{A = A}. No caso dos condicionais, a remoção de ramos só é feita quando não existe risco de eliminar labels que possam ser alvo de saltos.

\section{Geração de Código}

A geração de código foi implementada em \texttt{translator.py}. O tradutor percorre a AST e emite instruções da máquina virtual.

No início é gerada a instrução \texttt{START}; quando existem variáveis, é emitido \texttt{PUSHN} para reservar memória global; no fim é emitido \texttt{STOP}. O tradutor mantém uma tabela interna com o tipo, endereço, dimensão e indicação de array para cada variável.

\subsection{Expressões e Condições}

As expressões são traduzidas para a máquina de pilha. Por exemplo, assumindo \texttt{A} no endereço 0 e \texttt{B} no endereço 1:

\begin{lstlisting}
A = B + 1
\end{lstlisting}

gera código semelhante a:

\begin{lstlisting}
PUSHG 1
PUSHI 1
ADD
STOREG 0
\end{lstlisting}

Para valores reais são usadas instruções como \texttt{FADD}, \texttt{FSUB}, \texttt{FMUL} e \texttt{FDIV}. Condições usam instruções relacionais e saltos condicionais, como \texttt{JZ}.

Um exemplo simplificado de condicional é:

\begin{lstlisting}
IF (A .LT. B) THEN
    PRINT *, A
ENDIF
\end{lstlisting}

que é traduzido com uma comparação, um salto condicional para o fim do bloco e instruções de escrita.

\subsection{Ciclos DO e Arrays}

Os ciclos \texttt{DO} são traduzidos com labels internos e saltos. Para passos positivos, o ciclo continua enquanto a variável de controlo for menor ou igual ao limite; para passos negativos, continua enquanto for maior ou igual.

Os arrays são guardados como blocos contínuos na memória global. Como Fortran indexa a partir de 1, o endereço de \texttt{A(I)} é calculado como:

\begin{center}
endereço base + (índice - 1)
\end{center}

São usadas instruções como \texttt{PUSHGP}, \texttt{PADD}, \texttt{CHECK}, \texttt{LOAD} e \texttt{STORE}. A instrução \texttt{CHECK} valida limites em tempo de execução.

Um ciclo típico:

\begin{lstlisting}
DO 100 I = 1, N
    A(I) = A(I) + 1
100 CONTINUE
\end{lstlisting}

é traduzido usando uma label de início, uma comparação com o limite, um salto para o fim e uma atualização da variável de controlo. Esta parte foi uma das mais sensíveis, porque o label Fortran do \texttt{CONTINUE} tem de coincidir com o indicado na instrução \texttt{DO}.

\subsection{Subprogramas}

O compilador suporta subrotinas e funções simples. Como a máquina virtual não disponibiliza uma pilha de chamadas de alto nível para esta funcionalidade, os subprogramas são traduzidos diretamente no ponto de chamada. Por este motivo, chamadas recursivas não são suportadas.

Exemplo de subrotina:

\begin{lstlisting}
PROGRAM TESTSUB
INTEGER A, B
A = 14
B = 0
CALL SOMAUM(A, B)
PRINT *, A, B
END

SUBROUTINE SOMAUM(X, Y)
INTEGER X, Y
Y = X + 1
X = X + 1
RETURN
END
\end{lstlisting}

Para subrotinas, o tradutor usa aliases entre parâmetros formais e argumentos reais, permitindo que alterações em variáveis ou arrays sejam refletidas fora da subrotina.

As funções são tratadas de forma semelhante, mas devolvem um valor. Antes da execução do corpo da função, os argumentos são copiados para as posições de memória dos parâmetros formais. No fim, o valor da variável com o nome da função é colocado na pilha para ser usado pela expressão chamadora.

Esta estratégia funciona para os exemplos suportados, mas não implementa uma pilha de chamadas completa. Por isso, chamadas recursivas foram explicitamente excluídas.

\section{Testes}

Foram criados testes válidos e inválidos. Os testes válidos são compilados e comparados com o código VM esperado; os inválidos confirmam que erros sintáticos ou semânticos são detetados.

Entre os testes válidos encontram-se:

\begin{itemize}
    \item \texttt{somaarr.f77}: arrays e ciclos \texttt{DO};
    \item \texttt{primo.f77}: condições, ciclos e \texttt{MOD};
    \item \texttt{real\_logical.f77}: reais e lógicos;
    \item \texttt{do\_descendente.f77}: passo negativo;
    \item \texttt{subroutine\_scalar.f77}: subrotinas com escalares;
    \item \texttt{subroutine\_array.f77}: subrotinas com arrays;
    \item \texttt{conversor.f77}: funções em expressões;
    \item \texttt{optimization\_local.f77}: otimizações simples.
\end{itemize}

A execução automática é feita com:

\begin{lstlisting}
python3 tests/run_tests.py
\end{lstlisting}

Durante o desenvolvimento, estes testes foram úteis para encontrar problemas de geração de código. Em particular, os testes com arrays ajudaram a corrigir acessos deslocados por causa da indexação iniciada em 1, e os testes de subrotinas ajudaram a validar a passagem de escalares e arrays por referência.

\section{Dificuldades e Limitações}

As principais dificuldades foram:

\begin{itemize}
    \item cálculo correto de offsets em arrays, devido à indexação iniciada em 1;
    \item geração correta de labels e saltos para \texttt{DO}, \texttt{IF} e \texttt{GOTO};
    \item distinção entre chamada de função e acesso a array;
    \item passagem de argumentos para subrotinas, sobretudo arrays;
    \item compatibilização entre tipos inteiros, reais e lógicos.
\end{itemize}

As principais limitações atuais são:

\begin{itemize}
    \item apenas são suportados arrays unidimensionais;
    \item os subprogramas suportados cobrem apenas casos simples de \texttt{FUNCTION} e \texttt{SUBROUTINE};
    \item não há suporte a recursão;
    \item não existe uma pilha de chamadas completa, sendo os subprogramas traduzidos no ponto de chamada;
    \item o formato fixo clássico de Fortran 77 não é suportado integralmente;
    \item a recuperação de erros é simples.
\end{itemize}

Apesar destas limitações, o subconjunto implementado é suficiente para executar programas com cálculos inteiros e reais, manipulação de arrays, estruturas de controlo, leitura/escrita e subprogramas simples.

\section{Instruções de Utilização}

O compilador pode ser executado com:

\begin{lstlisting}
python3 compiler.py input.f77 output.vm
\end{lstlisting}

Exemplo:

\begin{lstlisting}
python3 compiler.py tests/somaarr.f77 tests/output/somaarr.vm
\end{lstlisting}

Requisitos:

\begin{itemize}
    \item Python 3;
    \item biblioteca PLY.
\end{itemize}

\section{Conclusão}

O projeto permitiu aplicar os principais conceitos abordados na unidade curricular: análise léxica, análise sintática, análise semântica, otimização simples e geração de código.

O compilador final processa programas com variáveis, arrays, ciclos, condições, leitura/escrita, labels, funções, subrotinas e expressões com diferentes tipos. Como trabalho futuro, seria interessante adicionar arrays multidimensionais, recursão, uma pilha de chamadas mais completa, melhor recuperação de erros e suporte mais fiel ao formato fixo clássico de Fortran 77.

\end{document}
